% ============================================================
%  Estado del arte y normativa
% ============================================================
\section{Estado del arte y normativas aplicables}
\label{sec:estado_arte}

\subsection{Panorama de los ensayos de seguridad eléctrica en la industria}

La seguridad eléctrica es un aspecto fundamental en la fabricación de productos
eléctricos y electrónicos, ya que garantiza tanto la integridad de los equipos como
la protección de los usuarios. Los ensayos de seguridad, que incluyen pruebas de
rigidez dieléctrica, continuidad de tierra y corriente de fuga, son requisitos
obligatorios para la certificación de productos en diversos mercados.
Tradicionalmente, estos ensayos se han realizado de manera manual, lo que puede dar
lugar a variabilidad en los resultados, prolongar los tiempos de prueba y generar
riesgos tanto para el personal como para los productos. Por esta razón, ha surgido
un creciente interés en la automatización de estos procesos, con el objetivo de
mejorar la precisión, la reproducibilidad y la eficiencia operativa, además de
garantizar la trazabilidad y un historial completo de las pruebas realizadas.
Existen en el mercado equipos comerciales que integran estos ensayos, como la
línea EFE Plus \cite{efeplus} o los equipos de Insatec \cite{insatec}, y
laboratorios que los ofrecen como servicio \cite{labconsultar}.

\subsection{Los ensayos de seguridad eléctrica más comunes}
\begin{itemize}
    \item \textbf{Rigidez dieléctrica (tensión resistida):} se aplica una tensión
          elevada entre las partes activas y las partes accesibles para verificar
          el aislamiento.
    \item \textbf{Continuidad de tierra:} verifica que la resistencia entre el
          borne de tierra y las partes metálicas accesibles sea baja.
    \item \textbf{Corriente de fuga:} asegura que las corrientes indeseadas hacia
          las partes accesibles estén dentro de los límites permitidos.
    \item \textbf{Consumo:} verifica que la potencia o la corriente absorbida no se
          aparte del valor nominal marcado más allá de la tolerancia admitida.
\end{itemize}

\subsection{Métodos tradicionales y su grado de automatización}

El proceso de prueba tradicional requiere intervención manual del operario y la
disposición del instrumental junto con personal capacitado, lo cual puede
producir variaciones en los resultados, demoras en las pruebas y posibles riesgos
para el personal, los instrumentos y los productos.

\subsection{Normativa aplicable}
\label{subsec:normativa}

La norma IEC~60335 establece requisitos de seguridad para aparatos eléctricos de
uso doméstico y aplicaciones comerciales de baja tensión. La parte~1
(IEC~60335-1) contiene los requisitos generales aplicables a todos los aparatos
comprendidos en la norma; en este trabajo se toma como referencia la
edición~4.1 \cite{iec60335}. La norma distingue los \textbf{ensayos de tipo}
(cláusulas~10 a~32), que certifican un diseño sobre muestras representativas, de
los \textbf{ensayos de rutina} (Anexo~A), que se aplican a cada unidad producida.
Por ser un equipo de final de línea, el sistema adopta los ensayos de rutina del
Anexo~A donde existen.

Otras normas de referencia contextual son:
\begin{itemize}
    \item \textbf{IEC~61010:} seguridad de equipos de medida, control y laboratorio
          \cite{iec61010}.
    \item \textbf{IEC~60950 / IEC~62368:} seguridad de equipos de tecnología de la
          información y audio/video \cite{iec62368}.
    \item \textbf{UL~60950 / UL~62368:} equivalentes para el mercado de EE.\,UU.
    \item \textbf{IRAM~2435:} norma argentina de referencia \cite{iram2435}.
    \item \textbf{ISO/IEC~17025:} competencia de laboratorios de ensayo y
          calibración \cite{iso17025}.
    \item \textbf{IEC~60990:} métodos de medición de la corriente de contacto y de
          la corriente del conductor de protección \cite{iec60990}.
\end{itemize}

\noindent Para el alcance de este proyecto se aplica la IEC~60335 y sus anexos.
La certificación del equipo ante el Organismo Argentino de Acreditación (OAA) u
otros organismos queda fuera del alcance, aunque el diseño se orienta al
cumplimiento normativo.
